\bigskip
\noindent Autres pistes de l'onglet Findings pour ce thème, dans des cotes qui n'ont pas de section ici :\par

\begin{encadre}
\begin{theoreme*}[Grothendieck, pages 1 à 3, groupoïde fondamental d'un produit fibré\piste{132-fibre-product-pi2-defect}]\nonverifie
Pour une fibration de Serre $f : X \to S$ et une application $g : Y \to S$, le foncteur $\Pi_1(X \times_S Y) \to \Pi_1(X) \times^2_{\Pi_1(S)} \Pi_1(Y)$ est bijectif sur $\pi_0$ et surjectif sur $\pi_1$, de noyau $\mathrm{Coker}(\pi_2X \times \pi_2Y \to \pi_2S)$ en la composante de base. C'est donc une équivalence quand $X$, $Y$ et $S$ sont asphériques.
\end{theoreme*}
\end{encadre}
{\small Le dernier mot de la page~3, \q{asphériques}, est illisible et restitué par la lecture. Mathlib n'a pas les fibrations de Serre ni le groupe $\pi_2$ à ce niveau. On se reportera aux notes de Grothendieck (\lot{132}{1}{1}).\par}

\begin{encadre}
\begin{conjecture*}[Grothendieck, pages 16 à 20, extensions d'un groupe non connexe ; voir aussi Brown--Mucuk 1994\piste{132-nonconnected-extension-classification}]
Soient $H$ un groupe topologique dont $\pi_0(H) = \mathfrak G$ est discret, et $\tilde H$ une extension de $H$ par $\pi = \pi_1(H^0)$ qui prolonge le revêtement universel de $H^0$. Pour tout groupe discret $N$, non nécessairement abélien, la catégorie des extensions de $H$ par $N$ est équivalente à celle des couples $(\tilde G_0, \varphi)$ formés d'une extension $\tilde G_0$ de $\mathfrak G$ par $N$ et d'un morphisme $\mathfrak G$-équivariant $\varphi : \pi \to \mathfrak z(N)$. L'extension $G$ s'obtient par produit contracté, et $\pi_0(G) \simeq \tilde G_0/\varphi(\pi)$.
\end{conjecture*}
\end{encadre}
{\small Les feuillets en esquissent une démonstration incomplète (pages 16 à 20) : seules les réductions et la construction explicite sont écrites, et les hypothèses de connexité locale sont de l'édition. Brown et Mucuk classent les revêtements des groupes topologiques non connexes par des modules croisés, non ces extensions. Mathlib n'a pas les revêtements universels des groupes topologiques. On se reportera aux notes de Grothendieck (\lot{132}{1}{16}).\par}

\begin{encadre}
\begin{exemple*}[Grothendieck, pages 5 à 8, Gr-catégorie périphérique d'une surface\piste{132-peripheral-gr-category}]\nonverifie
Pour une surface compacte orientable à $n \geq 1$ composantes de bord, le $2$-groupe des autoéquivalences du modèle formé de $\pi_1$, des lacets de bord $\ell_i$ et de leurs attaches $\alpha(i)$ a pour $\pi_0$ le groupe des classes d'applications à bord libre, et pour $\pi_1$ zéro, sauf pour le disque ($\mathbf Z$, canoniquement) et le cylindre ($\mathbf Z$, non canoniquement). Sa version à bord fixé a pour $\pi_0$ le groupe des classes d'applications relatif au bord, et un $\pi_1$ toujours nul.
\end{exemple*}
\end{encadre}
{\small La page calcule ces groupes sans les identifier aux groupes de classes d'applications. Cette identification, une forme de Dehn--Nielsen--Baer à plusieurs composantes de bord, est de la lecture, et n'est vérifiée que pour $n = 1$. Mathlib n'a pas les groupes de classes d'applications des surfaces. On se reportera aux notes de Grothendieck (\lot{132}{1}{5}).\par}

\begin{encadre}
\begin{conjecture*}[Grothendieck, pages 148 à 151, voisinage épointé et codimension deux\piste{147-deleted-neighbourhood-codim-two}]
Soient $X$ un espace analytique complexe lisse et $\Theta$ un diviseur à croisements normaux. Le voisinage tubulaire épointé $V^*_{\Theta,X}$ est la colimite homotopique du système semi-simplicial des voisinages tubulaires épointés des strates déployées. Son groupoïde fondamental est donc la somme amalgamée $\tilde\Pi_1 \leftarrow \tilde\Pi_{2,1} \to \tilde\Pi_2$ des groupoïdes des fibrés en tores multinormaux au-dessus des strates ouvertes de codimension~$1$ et~$2$. Il en va de même de $\Pi_1(X \setminus \Theta)$ quand $X \setminus \Theta$ et son bout ont des groupoïdes fondamentaux équivalents.
\end{conjecture*}
\end{encadre}
{\small La page~149 énonce le théorème de recollement comme une équivalence de topos et le corrige aussitôt en équivalence d'homotopie ; aucune démonstration n'est donnée. Mathlib n'a pas les voisinages tubulaires ni les colimites homotopiques d'espaces. On se reportera aux notes de Grothendieck (\lot{147}{8}{148}).\par}

\begin{encadre}
\begin{exemple*}[Grothendieck, pages 62 à 64, classe stratifiée de $M_{1,1}$ sur $\mathbf Z$\piste{147-m11-stratified-class-over-z}]\nonverifie
Dans l'invariant additif qui compte chaque pièce localement fermée d'un champ de Deligne--Mumford à groupe d'automorphismes constant $H$ comme la classe de son espace grossier divisée par $\mathrm{card}\,H$, le champ des courbes elliptiques sur $\operatorname{Spec}\mathbf Z$ a pour classe $\tfrac12(L-1) - \tfrac1{12} + \tfrac18\zeta_2 + \tfrac16\zeta_3$, où $\zeta_p = [\operatorname{Spec}\mathbf F_p]$, et sa compactification $\tfrac12 L - \tfrac1{12} + \tfrac18\zeta_2 + \tfrac16\zeta_3$. Sur $\mathbf Z[1/6]$, ces classes se réduisent à $\tfrac12(L-1) - \tfrac1{12}$ et $\tfrac12(L-1) + \tfrac5{12}$.
\end{exemple*}
\end{encadre}
{\small La page dit seulement \q{dans un groupe de Grothendieck convenable} ; la règle de l'invariant est reconstituée par la lecture, et l'arithmétique a été refaite. Mathlib n'a pas les champs de Deligne--Mumford. On se reportera aux notes de Grothendieck (\lot{147}{4}{62}).\par}

\begin{encadre}
\begin{conjecture*}[Grothendieck, page 18, l'ordre cyclique détermine la position\piste{154-cyclic-order-determines-position}]
Une position relative d'une pseudo-droite ajoutée, autre que les $D_i$ eux-mêmes, est déterminée par l'ordre cyclique, compatible à l'antipodie, dans lequel son revêtement double rencontre les pseudo-droites orientées de~$\Sigma$.
\end{conjecture*}
\end{encadre}
{\small La page écrit \q{il ne devrait pas être difficile de montrer} et ne le démontre pas. Mathlib n'a pas les arrangements de pseudo-droites. On se reportera aux notes de Grothendieck (\lot{154}{1}{18}).\par}

\begin{encadre}
\begin{theoreme*}[Grothendieck, pages 55 à 63, 85 et 86, monodromie diédrale\piste{154-dihedral-monodromy-half-turn}]\nonverifie
Soit $\Sigma$ un système de $n$ pseudo-droites d'un plan projectif réel. Les $2n$-gones $\mathrm{Pol}(D)$ découpés sur le revêtement double d'une pseudo-droite ajoutée $D$ forment un système local sur la surface cellulaire des positions relatives. Le transport le long d'un demi-tour autour d'une position $D = D_i$ est la rotation $t \mapsto t - (n-1)$ de $\mathbf Z/2n\mathbf Z$, et la monodromie $\pi_1 \to \mathbb D_{2n}$, à valeurs dans le groupe diédral d'ordre $4n$, est surjective sur la surface où la traversée des $D_i$ est permise.
\end{theoreme*}
\end{encadre}
{\small L'édition corrige deux fois $\mathbf Z/n\mathbf Z$ en $\mathbf Z/2n\mathbf Z$ (pages 57 et 85), et note que pour $n$ impair la surjectivité demande une réflexion latérale. La surjectivité sans traversée des $D_i$ (pages 85 et 86) n'est pas reprise. Mathlib n'a pas les arrangements de pseudo-droites ni les systèmes locaux sur un complexe cellulaire. On se reportera aux notes de Grothendieck (\lot{154}{3}{55}).\par}

\begin{encadre}
\begin{conjecture*}[Grothendieck, page 153, cube des générisations\piste{154-generisation-cube}]
Pour une position $D$ d'une pseudo-droite ajoutée passant par $\nu \geq 2$ sommets de $\Sigma$, les générisations stables de $D$ sont les $2^\nu$ sommets d'un cube de dimension $\nu$, dont les $\nu\,2^{\nu-1}$ arêtes sont les générisations par un seul sommet, l'incidence étant la spécialisation.
\end{conjecture*}
\end{encadre}
{\small La page l'énonce en mots, sans justification ; celle que donne la lecture, le côté où $D$ se décolle de chaque sommet se choisissant indépendamment, est la sienne. Une droite qui passe par $\nu$ sommets d'un arrangement de droites n'a que $2\nu$ perturbations stables, et pour $\nu \geq 3$ le compte est en tension avec la page~150. Mathlib n'a pas les arrangements de pseudo-droites. On se reportera aux notes de Grothendieck (\lot{154}{8}{153}).\par}

\begin{encadre}
\begin{conjecture*}[Grothendieck, pages 132 à 143 et 155 à 157, revêtement d'Arf par factorisation de Stein ; voir aussi Conrad~\cite{conrad2014}\piste{77-arf-cover-stein-factorisation}]
Sur toute base, caractéristique~$2$ comprise, les schémas de drapeaux isotropes $\Delta_H(E, q)$ d'un module quadratique régulier sont lisses et projectifs, à fibres géométriquement connexes sauf quand $N = 2n$ et $n \in H$. La factorisation de Stein de la grassmannienne isotrope maximale définit le revêtement double d'Arf, $\mathrm{Arf}' \simeq \mathrm{Arf} \wedge^{\mathbf Z/2} \mu_2$, et l'homomorphisme de Dickson est l'action de $O(q)$ sur ce revêtement, sans algèbres de Clifford.
\end{conjecture*}
\end{encadre}
{\small Les feuillets en esquissent une démonstration incomplète (pages 133 à 143) : la lissité et la connexité sont démontrées sur un corps, et le passage à une base quelconque \q{par spécialisation et factorisation de Stein} tient en une ligne. La section de la cote~77 ne traite pas cet énoncé. Conrad construit l'invariant de Dickson par l'algèbre de Clifford paire. Mathlib n'a pas les schémas de drapeaux isotropes ni la factorisation de Stein. On se reportera aux notes de Grothendieck (\lot{77}{7}{132}).\par}
